\bookchapter{Bab 63: View, State, dan Navigasi}

\chapterintro{Membuat mainan satu halaman sudah bisa. Bab ini menyelesaikan dua masalah “tumbuh dewasa”: bagaimana berpindah-pindah antar beberapa halaman (navigasi), dan bagaimana data mengalir antar halaman (manajemen state). Setelah menuntaskan bab ini, kamu bisa membangun kerangka App multi-halaman klasik seperti “daftar detail edit”.}

\booksection{Peta Besar Aliran Data dalam Manajemen State}

Alat state di SwiftUI tidak hanya @State. Lihat dulu peta besarnya untuk memahami “kapan memakai yang mana”:

\begin{center}
\begin{tabularx}{\linewidth}{|X|X|X|X|}
\hline
\textbf{Alat} & \textbf{Mengatur apa} & \textbf{Dipakai oleh} & \textbf{Skenario khas} \tabularnewline
\hline
@State & State sementara privat satu halaman & View saat ini & Penghitung, teks kolom input \tabularnewline
@Binding & Meminjam state milik view lain & View anak mengubah view induk & Komponen kustom \tabularnewline
@StateObject / @Observable & Logika yang dibagi lintas halaman & View model (ViewModel) & Pemuatan data, logika bisnis \tabularnewline
@EnvironmentObject / .environment & State yang dibagi secara global & Seluruh App & Status login pengguna, pengaturan tema \tabularnewline
\hline
\end{tabularx}
\end{center}

\begin{center}\small\color{gray}Peta alat state SwiftUI (pada 2026, @Observable sudah menjadi arus utama)\end{center}

\begin{lstlisting}[language=Swift]
// Cara modern di iOS 17+: makro @Observable (disarankan untuk proyek baru)
import Observation

@Observable
class CounterModel {                 // View model
    var count = 0
    func increment() { count += 1 }
}

struct CounterView: View {
    @State private var model = CounterModel()   // Memegang model

    var body: some View {
        VStack {
            Text("Kamu menekan \(model.count) kali")
            Button("Tekan aku") { model.increment() }
        }
    }
}
\end{lstlisting}

Saran blak-blakan untuk pemula: yang privat satu halaman pakai @State, yang lintas halaman/berisi logika pakai kelas @Observable. @StateObject/@ObservableObject di tutorial lama adalah gaya generasi sebelumnya — cukup pahami cara membacanya; untuk kode baru, ikuti panduan di atas.

\booksection{NavigationStack: “Tumpukan” untuk Pindah Halaman}

Sejak iOS 16, navigasi multi-halaman memakai NavigationStack — seperti tumpukan piring: satu masuk, satu ditaruh; saat kembali, dikeluarkan satu per satu:

\begin{lstlisting}[language=Swift]
struct HomeView: View {
    var body: some View {
        NavigationStack {
            List {
                NavigationLink("Buku kas Zhang San") {
                    DetailView(name: "Buku kas Zhang San")   // halaman tujuan
                }
                NavigationLink("Buku kas Li Si") {
                    DetailView(name: "Buku kas Li Si")
                }
            }
            .navigationTitle("Daftar buku kas")   // judul besar di atas
        }
    }
}

struct DetailView: View {
    let name: String
    var body: some View {
        Text("Halaman detail \(name)")
            .navigationTitle(name)
            .navigationBarTitleDisplayMode(.inline)
    }
}
\end{lstlisting}

NavigationLink adalah “tautan pindah halaman yang bisa diklik”; navigationTitle otomatis membuat judul bilah navigasi dan tombol kembali. Untuk pindah halaman secara terprogram (misalnya otomatis kembali setelah menekan “Simpan”), gunakan @Environment(\textbackslash{}.dismiss) untuk menutup halaman saat ini — dengan dua cara ini digabung, 90\% kebutuhan navigasi sudah tercakup.

\booksection{Popup dan Tab Bawah: sheet dan TabView}

Selain “mendorong halaman”, ada dua interaksi klasik lainnya:

\begin{lstlisting}[language=Swift]
// 1. Pop setengah layar/layar penuh dari bawah (favorit formulir input)
struct HomeView: View {
    @State private var showAdd = false
    var body: some View {
        Button("Catat") { showAdd = true }
            .sheet(isPresented: $showAdd) {
                AddRecordView()     // Halaman input yang muncul
            }
    }
}

// 2. Bar tab bawah ("kerangka" aplikasi)
struct MainTabView: View {
    var body: some View {
        TabView {
            RecordListView()
                .tabItem { Label("Buku kas", systemImage: "book") }
            StatsView()
                .tabItem { Label("Statistik", systemImage: "chart.pie") }
            SettingsView()
                .tabItem { Label("Pengaturan", systemImage: "gear") }
        }
    }
}
\end{lstlisting}

\begin{tcbinfo}{Tips: Kapan Memakai yang Mana?}
Aturan praktis: alur utama memakai NavigationStack untuk mendorong halaman (daftar detail), operasi sementara memakai popup sheet (tambah, edit, konfirmasi), dan beberapa bagian besar yang sederajat memakai TabView (tab bawah). Gabungan ketiganya adalah kerangka dari hampir semua App.

\end{tcbinfo}

\booksection{List: “Daftar Andalan” iOS}

Komponen paling sering dipakai sekaligus paling khas Apple di iOS adalah List — garis pemisah, hapus dengan geser, dan sorotan saat diklik sudah tersedia secara bawaan:

\begin{lstlisting}[language=Swift]
struct RecordListView: View {
    @State private var records = ["Makan siang 25.5", "MRT 6.0", "Gaji +8000"]

    var body: some View {
        List {
            ForEach(records, id: \.self) { r in
                Text(r)
            }
            .onDelete { indexSet in          // geser kiri untuk menghapus
                records.remove(atOffsets: indexSet)
            }
        }
    }
}
\end{lstlisting}

Daftar + NavigationLink + sheet + TabView — gabungkan empat hal ini dan kamu sudah bisa membangun kerangka App lengkap “daftar catatan buka detail tekan + untuk entri”. Proyek buku catatan keuangan di bab berikutnya dibangun persis seperti ini.

\begin{tcbpractice}{Praktik: Membangun "Kerangka Tiga Halaman"}
\booksection{Siklus Hidup Aplikasi: Latar Depan dan Latar Belakang}

Mirip siklus hidup Activity di Android, App iOS juga mengalami perubahan status “lahir, latar depan, latar belakang, mati”. Di SwiftUI, amati lewat nilai lingkungan scenePhase:

\begin{lstlisting}[language=Swift]
@main
struct MyApp: App {
    @Environment(\.scenePhase) private var phase

    var body: some Scene {
        WindowGroup {
            ContentView()
        }
        .onChange(of: phase) { _, newPhase in
            switch newPhase {
            case .active:    print("Kembali ke depan, lanjutkan pembaruan")
            case .inactive:  print("Sedang transisi (panggilan masuk/beralih ke latar sesaat)")
            case .background: print("Masuk latar, simpan draf")
            @unknown default: break
            }
        }
    }
}
\end{lstlisting}

Kegunaan khas: menyegarkan data saat kembali ke latar depan, menyimpan draf saat masuk latar belakang. iOS sangat ketat soal berjalan di latar belakang — App biasa segera ditangguhkan (dibekukan) tidak lama setelah masuk latar belakang, jadi jangan berharap bisa diam-diam mengerjakan tugas berat di sana; untuk tugas terjadwal, gunakan notifikasi lokal dari Bab 68, bukan “tetap hidup di latar belakang”.

Buat kerangka buku catatan keuangan versi data palsu: TabView dengan tiga tab (buku catatan/statistik/pengaturan) halaman buku catatan memakai List untuk menampilkan data palsu, ketuk baris untuk ke detail (NavigationStack) tombol “+” di pojok kanan atas memunculkan sheet entri. Tanpa menyimpan data, tanpa internet, murni latihan kerangka — setelah kerangka ini jadi, mengisi data sungguhan di Bab 65 akan berjalan mulus.

\end{tcbpractice}

\begin{tcbsummary}{Ringkasan Bab Ini}
\begin{itemize}
  \item Empat alat state: @State (privat satu halaman), @Binding (pinjam), @Observable (logika), .environment (global)
  \item Proyek baru tulis view model dengan makro @Observable; @ObservableObject di tutorial lama cukup dibaca saja
  \item NavigationStack mendorong halaman; @Environment(\textbackslash{}.dismiss) untuk kembali secara terprogram
  \item Alur utama mendorong halaman, operasi sementara memunculkan sheet, bagian besar yang sederajat memakai TabView
  \item List adalah "daftar andalan": garis pemisah dan hapus dengan geser langsung siap pakai
\end{itemize}

\end{tcbsummary}
